Der Tropfen hat die Schmelztemperatur $\vartheta_\text{f} = \TfPbO$. Im Wasser erstarrt er (Erstarrungswärme $L_\text{f} = \LfPbO$) und kühlt dann als festes Blei ($c_\text{Pb} = \cPbO$, FoTaBe S.~190) auf die Endtemperatur $\vartheta_1+\Delta\vartheta$ ab. Das Wasser mit $m_\text{W} = \rho\,V = \mWa$ erwärmt sich von $\vartheta_1$ um $\Delta\vartheta$:
\begin{align}
 m_\text{Pb}\left(L_\text{f}+c_\text{Pb}\left(\vartheta_\text{f}-\vartheta_1-\Delta\vartheta\right)\right) &= m_\text{W}c_\text{W}\,\Delta\vartheta
\end{align}
Aufgelöst nach der Anfangstemperatur:
\begin{align}
 \vartheta_1 &= \TaF \\
   &= \TfPbO-\dT-\frac{1}{\cPb}\cdot\left(\frac{\mWa\cdot\cW\cdot\dT}{\mPb}-\LfPb\right) \\
   &= \Ta \approx \result{\TaP{0}{2}}
\end{align}
Das Wasser war also schon recht warm. Weil der Temperaturanstieg nur auf ein Zehntelgrad genau angegeben ist, ist das Ergebnis nur grob.
